% SOURCE_AUTHORITY: sga5_fr_workpass.tex, lines 9903--9945 (LF-normalized unit).
% SOURCE_SHA256: 791F4EFFC5E02832D5D77ED03518C8156D6F07E4C8238B03545DB93D883FBB28
\begin{statement}{Corolario 4.6}
Sean $S$ un esquema, $E$ un $\mathcal O_S$-módulo localmente libre de rango constante $r$ y $g:X=\V(\check E)\to S$ el fibrado vectorial cuyo haz de secciones es $E$. Para toda sección $f$ de $X$ sobre $S$, se tiene
\[
f_*(1_S)=g^*(c_r(E))
\]
en $H^{2r}(X,\mu_X^{\otimes r})$.
\end{statement}

\emph{Demostración.} Como $g^*$ es un isomorfismo (1.2) y $g\circ f=\mathrm{id}$, $f^*$ es biyectivo. Basta, por tanto, ver que
\[
f^*f_*(1_S)=f^*g^*(c_r(E)).
\]
Ahora bien, denotando por $N$ el haz normal de $f$, esta última igualdad equivale, por 4.1, a
\[
c_r(\check N)=c_r(E).
\]
De hecho, se tiene incluso un isomorfismo $\check N\simeq E$, pues la sucesión exacta de (EGA IV 17.2.5), aplicada a la situación
\[
S\xrightarrow{f} X\xrightarrow{g} S,
\]
proporciona un isomorfismo $N\simeq f^*\Omega^1_{X/S}$ y, por otra parte, $\Omega^1_{X/S}\simeq g^*\check E$ (EGA IV 16.4.8).

\begin{statement}{Corolario 4.7}
Bajo las hipótesis de 4.6, se supone además que $S$ es liso y cuasiproyectivo sobre un cuerpo separablemente cerrado. Entonces, denotando por $\sigma_X$ el ciclo de la sección nula de $X$ y por $Y=f^{-1}(\sigma_X)$ el ciclo de ceros de $f$, se tiene
\[
\cl_S(Y)=c_r(E)
\]
en $H^{2r}(S,\mu_S^{\otimes r})$.
\end{statement}

\emph{Demostración.} Por IV,
\[
\cl_S(Y)=f^*\cl_X(\sigma_X).
\]
Denotando por $s_0$ el morfismo de sección nula de $X$, se tiene
\[
\cl_X(\sigma_X)\stackrel{\mathrm{dfn}}{=}(s_0)_*(1_S)\stackrel{(4.6)}{=}g^*(c_r(E)),
\]
de donde se sigue la afirmación aplicando el homomorfismo $f^*$ a los términos extremos.

\begin{statement}{Observación 4.8}
Sea $\ell$ un número primo distinto de las características residuales de $S$. Por un simple paso al límite proyectivo, se deducen de 4.6 y 4.7 los enunciados correspondientes en cohomología $\ell$-ádica.
\end{statement}
